Dans un pentagone convexe, on peut mesurer la somme de ses angles
intérieurs. En traçant plusieurs pentagones à main levée, on trouve souvent
approximativement $540^\circ$.
\begin{itemize}
  \item Expliquez pourquoi on ne peut pas se contenter d'exemples tracés et
  mesurés pour conclure.
  \item Donnez un argument géométrique \emph{global} prouvant que la valeur de
  la somme des angles d'un pentagone donne toujours $540^\circ$ (indice :
  triangulez la figure).
\end{itemize}
